\specialsection{Введение}

Севооборот относится к числу базовых практик земледелия, поскольку последовательность культур на одном и том же поле
влияет на дальнейшее использование участка и на характер принимаемых агрономических решений. Выбор следующей культуры
обычно не является случайным: он связан с предшествующей историей поля, особенностями региона и накопленными
закономерностями смены культур. Поэтому анализ таких последовательностей представляет интерес как с практической, так и
с исследовательской точки зрения.

Развитие цифровых источников данных в сельском хозяйстве делает возможным переход от ручного анализа к более
формализованным методам обработки исторической информации по полям. Если ранее решение о следующей культуре в основном
принималось на основе опыта специалиста и локальных наблюдений, то наличие накопленных последовательностей культур
позволяет использовать методы анализа данных и машинного обучения для поиска устойчивых паттернов и формирования
рекомендаций. При этом на практике важно не только получить один наиболее вероятный вариант, но и рассмотреть несколько
правдоподобных альтернатив.

Наряду с этим для систем подобного рода важна интерпретируемость результата. В прикладном сценарии пользователю
недостаточно знать только рекомендуемую культуру; важно также понимать, насколько уверенно сделана рекомендация и какие
факторы поддерживают или ослабляют отдельных кандидатов. Поэтому при разработке подобных систем актуальной становится не
только задача прогнозирования, но и задача объяснения полученного результата.

В данной работе рассматривается построение исследовательского прототипа рекомендательной системы, которая по истории
культур на конкретном поле и базовым признакам поля формирует ранжированный список вероятных следующих культур. При этом
работа сознательно не претендует на создание полноценного агрономического оптимизатора и не рассматривает систему как
замену экспертному решению. Основной акцент делается на построении воспроизводимого и интерпретируемого прототипа,
который может использоваться как инструмент поддержки анализа.

Работа сочетает методы анализа последовательностей культур, многоклассового прогнозирования, рекомендательных систем
и интерпретации результатов. Это позволяет рассматривать задачу выбора следующей культуры не только как задачу предсказания,
но и как задачу формирования интерпретируемого списка вероятных кандидатов для поддержки принятия решений.

В первой части работы рассматриваются постановка задачи, её границы и место работы среди существующих подходов. Далее
описываются общая архитектура решения, используемые данные и этапы их подготовки. После этого анализируются ядро
модели, рекомендательный режим, оценка уверенности, слой интерпретируемости и пространственная демонстрация результатов.
